
% !TEX program=pdflatex % Usa il compilatore «PDFLaTeX»
% LTeX: language=en-GB

\documentclass[
    conference,
    nofonttune,
    11pt
]{IEEEtran}
% \IEEEoverridecommandlockouts
% The preceding line is only needed to identify funding in the first footnote. If that is unneeded, please comment it out.


\makeatletter%
\newcommand\checkClass[2]{%
    \@ifclassloaded{IEEEtran}{#1}{#2}%
}%
\makeatother%

% \checkClass{% «IEEEtran» class
% }{% «Beamer» class
% }%

\usepackage[british]{babel}
\usepackage{microtype}
\usepackage{fontspec}

\checkClass{% «IEEEtran» class
    \setmainfont{cambria}[%
        Path = ../References/, % Pathc
        UprightFont = cambria.ttc,    % Normal font
        ItalicFont = cambriai.ttf,    % Italic font
        BoldFont = cambriab.ttf,      % Bold font
        BoldItalicFont = cambriaz.ttf % Italic bold font
    ]%
}{% «Beamer» class
    \setsansfont{cambria}[%
        Path = ../References/, % Pathc
        UprightFont = cambria.ttc,    % Normal font
        ItalicFont = cambriai.ttf,    % Italic font
        BoldFont = cambriab.ttf,      % Bold font
        BoldItalicFont = cambriaz.ttf % Italic bold font
    ]%
}%

\usepackage{cite}
\usepackage{amsmath,
            amsfonts,
            amsthm,
            amssymb}

\usepackage{mathtools}
\usepackage{algorithmic}
\usepackage{textcomp}


\usepackage{graphicx}
\usepackage{xcolor}

\usepackage{tikz}
% \usetikzlibrary{
%     tikzmark,
%     positioning,
%     arrows.meta,
%     fit,
%     calc,
%     decorations.pathreplacing,
%     calligraphy,
%     spy,
% }
\usepackage{pgfplots}
\usepgfplotslibrary{groupplots}
\pgfplotsset{compat=1.18}

\usepackage{showframe}
\renewcommand\ShowFrameLinethickness{.5pt}    % line thickness
\renewcommand\ShowFrameColor{\color{magenta}} % line color


\usepackage[shortlabels]{enumitem}
\setlist[itemize]{label=$\diamond$}

% \usepackage{caption,
%             subcaption}
\usepackage{float,
            stfloats}

\usepackage{hyperref}
\newcommand{\teqref}[2]{\hyperref[#1]{#2}}


\usepackage{cleveref}
%
\crefrangeformat{equation}{(#3#1#4--#5#2#6)}%
\crefmultiformat{equation}{(#2#1#3}{, #2#1#3)}{, #2#1#3}{, #2#1#3)}%
\crefrangemultiformat{equation}{(#3#1#4--#5#2#6}{, #3#1#4--#5#2#6)}{, #3#1#4--#5#2#6}{, #3#1#4--#5#2#6)}%
%
\crefrangeformat{enumi}{#3#1#4--#5#2#6}%
\crefmultiformat{enumi}{#2#1#3}{ and #2#1#3}{, #2#1#3}{ and #2#1#3}%
%
\crefrangeformat{figure}{#3#1#4--#5#2#6}%
\crefmultiformat{figure}{#2#1#3}{ and #2#1#3}{, #2#1#3}{ and #2#1#3}%
\crefrangemultiformat{figure}{#3#1#4--#5#2#6}{ and #3#1#4--#5#2#6}{, #3#1#4--#5#2#6}{ and #3#1#4--#5#2#6}%
%
\crefrangeformat{section}{#3#1#4--#5#2#6}%
\crefmultiformat{section}{#2#1#3}{ and #2#1#3}{, #2#1#3}{ and #2#1#3}%


\usepackage{siunitx}
\sisetup{
    inter-unit-product=\ensuremath{{}\cdot{}},
    scientific-notation=false,
}

\usepackage{comment}

\def\BibTeX{%
    {\rm B\kern-.05em{\sc i\kern-.025em b}\kern-.08em%
    T\kern-.1667em\lower.7ex\hbox{E}\kern-.125emX}%
}%

\newcommand\dimTesto[1]{%
    \pgfmathparse{#1*1.2}%
    \fontsize{#1 pt}{\pgfmathresult pt}\selectfont%
}

\newcommand\derT[3][d]{\frac{\mathrm{#1}#2}{\mathrm{#1}#3}} % Derivata totale
\newcommand\derTiL[3][d]{\mathrm{#1}#2/\mathrm{#1}#3}       % Derivata totale in linea
\newcommand\derP[2]{\frac{\partial#1}{\partial#2}}	        % Derivata parziale

\newcommand\de{\mathrm{d}}

\newcommand{\SpazMate}[3]{% Spaz[iatura] Mate[matica]
    \thinmuskip=#1mu\relax
    \medmuskip=#2mu\relax
    \thickmuskip=#3mu\relax
}
\SpazMate{1}{2}{3} % Seleceting math spacing

\newcommand\ie{\textit{i.e.}}

\newcommand\vColla[3]{\vspace{#1#3 plus #2#3 minus #2#3}}
\newcommand\hColla[3]{\hspace{#1#3 plus #2#3 minus #2#3}}

\title{
    \vspace*{-1cm}
    \textbf{Analysis of different MOR methods for solving nonlinear Elliptic PDEs}
}

\author{
    \IEEEauthorblockN{1\textsuperscript{st} Maddalena Ghiotti}
    \IEEEauthorblockA{s332834}
    \and
    \IEEEauthorblockN{2\textsuperscript{nd} Valerio Taralli}
    \IEEEauthorblockA{s332689}
    \and
    \IEEEauthorblockN{3\textsuperscript{rd} Gaia Saroglia}
    \IEEEauthorblockA{s332687}
}

\begin{document}
    
    \twocolumn[
        \begin{@twocolumnfalse}
            \begin{center}
                \noindent\parbox{.75\textwidth}{\maketitle}
                \noindent\parbox{.9\textwidth}{%
                \vspace{-.5cm}
                \begin{abstract}
                    \bfseries
                    \noindent PUNTI DA SCRIVERE (cenni): 
                    \begin{itemize}
                        \item cosa facciamo e di cosa si tratta
                        \item descrizione problema (su cosa facciamo esperimenti numerici): problema ellittico non lineare parametrizzato, condizioni al contorno, configurazioni in analisi  
                        \item obiettivo: valutare e confrontare POD e PINNs+POD-NN rispetto a full order 
                        \item risultati: cosa abbiamo ottenuto dalle simulazioni
                    \end{itemize}
                \end{abstract}}
            \end{center}
            % Per qualche motivazione «abstract» richiede «\text{}» per le parentesi quadre
        \end{@twocolumnfalse}
        \vspace{0.5cm}
    ]
    
    % \begin{comment}
    % \begin{IEEEkeywords}
    % component, formatting, style, styling, insert
    % \end{IEEEkeywords}
    % \end{comment}

    %
    
    \section{Introduction}\label{secIntroduction}

        The numerical approximation of nonlinear parametrized partial differential equations (PDEs) using high fidelity techniques, which means solving the full discretized problem, called full-order model (FOM), typically results in very large systems that become computationally expensive and often unfeasible. For this reason, model order reduction (MOR) techniques have gained considerable attention: they are exploited to decrease computational costs while maintaining a good solution approximation.
        
        In the present paper, we study a nonlinear elliptic parametrized equation with homogeneous Dirichlet boundary conditions, on a two-dimensional square domain: both cases with source term dependent on or independent of the parameter are considered. In these two settings, three approaches are implemented\footnote{For the acronyms meaning see later sactions.}:

        \begin{enumerate}%[leftmargin=.5cm]
            \item the first is the POD-Galerkin method over a finite element FOM;
            \item then we focus on parametric PINN, using the same network structure with either parameters independent and dependent source term;
            \item in the end, we use, as another reduce method, the hybrid approach of the POD-NN.
        \end{enumerate}

        Afterwards, the results in terms of computational costs and accuracy are compared using as a benchmark the FOM results. 
        
        The outline of the project is as follows: we first introduce the model definition and the methods used; the numerical studies, results and comparisons are next tackled for the given problems; lastly, the conclusions are discussed.

        Before continuing, we assume the reader is proficient enough about functional analysis, to avoid explaining all the Hilbert spaces used in the weak formulation (such as $H_0^1(\Omega)$), and about FEM.

    %
    
    \section{PDE definition and methods}\label{secModelDef}
    
        In the following, we talk about the discussed PDE and briefly describe the four approaches studied to solve it.

        It's noteworthy that, for the sake of notational brevity, the [explicit] dependency on the parameters $\boldmu$ in any function $f$ is represented as a superscript: $f^\boldmu\equiv f(\mBF x;\boldmu)$.
        % Explicit dependency means that a function depend on the parameters regardless of its arguments like (k^μ(·,·) with the ratio μ₀/μ₁ and the exponent exp(μ₁u^μ)); conversely, if the function depends on the parameters only through one of its arguments, then it should not have a superscript (like a(·,·) with u^μ)

        \subsection{PDE definition}

            The nonlinear elliptic PDE studied consist of finding
            $$
                u^\boldmu(\mBF x):\Omega\times\mathcal P\to\mBB R
            $$
            such that it's the solution, for some fixed parameters
            $$
                \boldmu\equiv(\mu_0,\mu_1)\in\mathcal P=[0.1,1]^2
            $$
            in the parametric square domain, of the following problem with homogeneous Dirichlet boundary conditions:
            %
            \begin{equation}
                \refstepcounter{equation}
                \label{eqNLEllipticPDE}
                \newcommand\spzOriz{\hspace{.5cm}}
                \text{\NLEllPDE}
                \left\{\begin{alignedat}{2}
                    &-\Delta u^\boldmu
                    +h^\boldmu
                    =g^\boldmu+\frac{\mu_0}{\mu_1}&\spzOriz&
                    \text{in }\Omega,\\
                    &u^\boldmu=0&\spzOriz&
                    \text{on }\partial\Omega,
                \end{alignedat}\right.
                \notag
            \end{equation}
            %
            where
            %
            \begin{equation*}
                % \label{eqNonLinearTerm}
                h^\boldmu:V\times\mathcal P\to V,\quad
                (u;\boldmu)\mapsto
                h^\boldmu(u)\equiv
                \frac{\mu_0}{\mu_1}e^{\mu_1 u}
            \end{equation*}
            %
            is the nonlinear functional term on a space $V$ later defined,
            $$
                g^\boldmu(\mBF x):\Omega\times\mathcal P\to\mBB R
            $$
            is the source term to be defined while
            $$
                \mBF x\equiv(x_0,x_1)\in\Omega=(0,1)^2
            $$
            is the spatial square domain. For the source term two different cases are considered:
            
            \begin{enumerate}[
                label=C\arabic*.,
                rightmargin=0cm
            ]
                \item\label{hypIndependentSourceTerm} independent of $\boldmu$ with $g^\boldmu(\mBF x)\equiv g_1(\mBF x)$ where:
                    %
                    \begin{equation}
                        \label{eqIndependentSourceTerm}
                        \begin{aligned}
                            g_1(\mBF x)&\equiv
                            100\sin(2\pi x_0)\cos(2\pi x_1)\\[-3.5ex]
                            &\hphantom{\equiv100\sin(2\pi\mu_0x_0)\cos(2\mu_0\pi x_1)}
                        \end{aligned}
                        %   \quad\forall\mBF x\in\Omega
                    \end{equation}
            
                \item\label{hypDependentSourceTerm} dependent on $\boldmu$ with $g^\boldmu(\mBF x)\equiv g^{\mu_0}_2(\mBF x)$ where:
                    %
                    \begin{equation}
                        \label{eqDependentSourceTerm}
                        g^{\mu_0}_2(\mBF x)\equiv
                        100\sin(2\pi\mu_0x_0)\cos(2\mu_0\pi x_1)
                        %   \quad\forall\mBF x\in\Omega
                    \end{equation}
            \end{enumerate}
            
            Lastly the weak formulation of \refNLEllPDE{} is
            %
            \begin{equation}
                \refstepcounter{equation}
                \label{eqVarNLEllipticPED}
                % \vcenter{\hbox{\rotatebox{90}{\text{\VarNLEllPDE}}}}
                \text{\VarNLEllPDE}
                \left\{\begin{aligned}
                    &\text{Find $u^\boldmu\in V$
                    such that $\forall v\in V$}\\
                    &a(u^\boldmu,v)
                    +k^\boldmu(u^\boldmu,v)=
                    F^\boldmu(v),
                \end{aligned}\right.
                \notag
            \end{equation}
            %
            where $V\equiv H_0^1(\Omega)$
            %
            \begin{equation*}
                % \SpazMate{1}{1}{1}
                \newcommand\sepSpace{\quad}
                \begin{alignedat}{2}
                    a:&V\times V\to\mBB R,%\\&&&
                    &\sepSpace(u,v)&\mapsto a(u,v)\equiv
                    \int_\Omega\nabla u\cdot\nabla v;\\
                    k^\boldmu:&V\times V\to\mBB R,%\\&&&
                    &\sepSpace(u,v)&\mapsto k^\boldmu(u,v)\equiv
                    \int_\Omega h^\boldmu(u)v;\\
                    F^\boldmu:&V\to\mBB R,
                    &\sepSpace v&\mapsto F^\boldmu(v)\equiv
                    \int_\Omega\left(g^\boldmu+\frac{\mu_0}{\mu_1}\right)v,
                \end{alignedat}
            \end{equation*}
            %
            for which we assume that $a$ is continuous and coercive with constants $\gamma\in\mBB R<+\infty$ and $\alpha\in\mBB R>0$ respectively.
            
        \subsection{Full order model (FOM)}

            The FOM solution uses the Galerkin method: consider a subspace $V_\delta\subseteq V$ with dimension $N_\delta\equiv\dim\{V_\delta\}_{\delta>0}$, depending on the refinement parameter $\delta>0$, then \refVarNLEllPDE{} becomes
            %
            \begin{equation}
                \refstepcounter{equation}
                \label{eqGalerVarNLEllipticPED}
                % \vcenter{\hbox{\rotatebox{90}{\text{\VarGalerNLEllPDE}}}}
                \text{\VarGalerNLEllPDE}
                \left\{\begin{aligned}
                    &\text{Find $u^\boldmu_\delta(\boldmu)\in V_\delta$
                    such that $\forall v_\delta\in V_\delta$}\\
                    &a(u^\boldmu_\delta,v_\delta)
                    +k^\boldmu(u^\boldmu_\delta,v_\delta)=
                    F^\boldmu(v_\delta),
                \end{aligned}\right.
                \notag
            \end{equation}
            %
            in which now $u^\boldmu_\delta$ approximate the true solution $u^\boldmu$. However, due to the nonlinearity of $k^\boldmu$ it's not possible to obtain a matrix form of \refVarGalerNLEllPDE{}; hence, it's necessary to use an iterative approach, chosen to be the Newton method: assume to know the solution of $u^\boldmu_{\delta,n}$ of the $k$-th step, the update rule takes the form
            %
            \begin{equation}
                \label{eqUpdateRuleNewtonMethod}
                u^\boldmu_{\delta,n+1}
                =u^\boldmu_{\delta,n}
                +\delta u^\boldmu_n.
                % \overset{k\to\infty}\longrightarrow.
            \end{equation}
            %
            Given the residual of \refVarGalerNLEllPDE{}
            $$
                r^\boldmu(u^\boldmu_\delta,v_\delta)\equiv
                a(u^\boldmu_\delta,v_\delta)
                +k^\boldmu(u^\boldmu_\delta,v_\delta)
                -F^\boldmu(v_\delta)=0,
            $$
            the direction function $\delta u$ can be found by having $u^\boldmu_{\delta,n+1}$ such that $r^\boldmu(u^\boldmu_{\delta,n+1},v)=r^\boldmu(u^\boldmu_{\delta,n}+\delta u^\boldmu_n,v)=0$, which when linearized gives
            $$
                0=r^\boldmu(u^\boldmu_{\delta,n}+\delta u,v_\delta)
                \approx r^\boldmu(u^\boldmu_{\delta,n},v_\delta)
                +J^\boldmu_{r,n}[\delta u^\boldmu_n,v_\delta]\delta u^\boldmu_n,
            $$
            that is
            %
            \begin{equation}
                \label{eqDeltauConstraint}
                J^\boldmu_{r,n}[\delta u^\boldmu_n,v_\delta]\delta u^\boldmu_n
                =-r^\boldmu(u^\boldmu_{\delta,n},v_\delta)
                \quad\forall v_\delta\in V_\delta,
            \end{equation}
            %
            in which $J^\boldmu_{r,n}[\delta u^\boldmu_n,v_\delta]$ is the Jacobian of the residual $r^\boldmu$ evaluated in $u^\boldmu_{\delta,n}$ along the unknown direction $\delta u^\boldmu_n$, namely
            %
            \begin{equation}
                \label{eqGateauxDerivativeDefinition}
                J^\boldmu_{r,n}[\delta u^\boldmu_n,v_\delta]\equiv\lim_{h\rightarrow0}
                \frac{r^\boldmu(u^\boldmu_{\delta,n}+h\delta u^\boldmu_n,v_\delta)
                -r^\boldmu(u^\boldmu_{\delta,n},v_\delta)}h,
            \end{equation}
            %
            assuming $r^\boldmu$ is Gateaux differentiable. From \eqref{eqGateauxDerivativeDefinition}, it's then clear that
            $$
                J_{a,n}[\delta u^\boldmu_n,v_\delta]
                =\int_\Omega\nabla\delta u^\boldmu_n\nabla v_\delta,
            $$
            for the linearity of $a(u^\boldmu_\delta,v_\delta)$ w.r.t. $u^\boldmu_\delta$, whereas
            $$
                J^\boldmu_{F,n}[\delta u^\boldmu_n;v_\delta]=0
            $$
            as $F^\boldmu$ it does not depend on $u^\boldmu_\delta$; thus, the only Jacobian necessary to calculate explicitly is
            %
            \begin{equation*}
                % \SpazMate{1}{1}{1}
                \begin{aligned}
                    J^\boldmu_{k,n}[\delta u^\boldmu_n,v_\delta]
                    % &=\lim_{h\rightarrow0}{\frac1h\left[\int_\Omega
                    % \frac{\mu_0}{\mu_1}e^{\mu_1(u^\boldmu_{\delta,n}+h\delta u)}v_\delta
                    % -\int_\Omega\frac{\mu_0}{\mu_1}e^{\mu_1u^\boldmu_{\delta,n}}v_\delta\right]}
                    % \\&=\lim_{h\rightarrow0}{\frac1h\left[\int_\Omega
                    % \frac{\mu_0}{\mu_1}\left(e^{\mu_1u^\boldmu_{\delta,n}}e^{\mu_1h\delta u}
                    % -e^{\mu_1u^\boldmu_{\delta,n}}\right)v_\delta\right]}\\
                    &=\lim_{h\rightarrow0}{\frac1h\left[\int_\Omega\frac{\mu_0}{\mu_1}
                    \left(e^{\mu_1h\delta u}-1\right)e^{\mu_1u^\boldmu_{\delta,n}}v_\delta\right]}
                    \\&=\lim_{h\rightarrow0}{\frac1h\left[\int_\Omega\frac{\mu_0}{\mu_1}
                    \left(1+\mu_1h\delta u-1\right)e^{\mu_1u^\boldmu_{\delta,n}}v_\delta\right]}
                    \\&=\int_\Omega\mu_0\delta ue^{\mu_1u^\boldmu_{\delta,n}}v_\delta,
                \end{aligned}
            \end{equation*}
            %
            where $e^{\mu_1h\delta u}\approx1+\mu_1h\delta u$ is valid only with $h\ll1$ as is the case when $h\rightarrow0$. Therefore, \eqref{eqDeltauConstraint} becomes
            %
            \begin{equation}
                \refstepcounter{equation}
                \label{eqIterativeVarNLEllipticPED}
                % \vcenter{\hbox{\rotatebox{90}{\text{\NewtonNLEllPDE}}}}
                \text{\NewtonNLEllPDE}
                \left\{\begin{aligned}
                    &\text{Find $\delta u^\boldmu_n\in V_\delta$
                    such that $\forall v_\delta\in V_\delta$}\\
                    &\begin{aligned}
                        &a(\delta u^\boldmu_n,v_\delta)
                        +\mu_0b(\delta u^\boldmu_n,
                        v_\delta e^{\mu_1u^\boldmu_{\delta,n}})
                        =\\&\quad-a(u^\boldmu_{\delta,n},v_\delta)
                        -\frac{\mu_0}{\mu_1}
                        b(e^{\mu_1u^\boldmu_{\delta,n}},v_\delta)
                        +F^\boldmu(v_\delta),
                    \end{aligned}
                \end{aligned}\right.
                \notag
            \end{equation}
            in which
            %
            \begin{equation*}
                b:V\times V\to\mBB R,\quad
                (u,v)\mapsto b(u,v)
                \equiv\int_\Omega u v.
            \end{equation*}
            %
            Introducing now a set of basis function $\{\varphi_k\}_{k=1}^\Ndelta$ such that $V_\delta=\spanOp\{\varphi_k\}_{k=1}^\Ndelta$, leads to the linear system
            %
            \begin{equation}
                \label{eqGalerIterLS}
                \mBF S^\boldmu_n\delta\mBF u^\boldmu_n
                =\mBF y^\boldmu_n,
            \end{equation}
            %
            where
            $$
                \mBF S^\boldmu_n
                \equiv\mBF S^\boldmu(\mBF u^\boldmu_{\delta,n})
                \in\mBB R^{\Ndelta\times\Ndelta}
                \quad\text{and}\quad
                \mBF y^\boldmu_n
                \equiv\mBF y^\boldmu(\mBF u^\boldmu_{\delta,n})
                \in\mBB R^\Ndelta
            $$
            have the following definition
            %
            \begin{equation*}
                \begin{aligned}
                    \mBF S^\boldmu_n&\equiv
                    \mBF A+\mu_0\mBF B^{\mu_1}_n,\\
                    \mBF y^\boldmu_n&\equiv-\mBF a_n
                    -\frac{\mu_0}{\mu_1}(\mBF b^{\mu_1}_n-\mBF b)
                    +\mBF g^\boldmu.
                \end{aligned}
            \end{equation*}
            %
            Some remarks ensue from \eqref{eqGalerIterLS}:

            \begin{itemize}[
                itemsep=2pt,
                % rightmargin=0cm
            ]
                \item The update rule in \eqref{eqUpdateRuleNewtonMethod} becomes
                    %
                    \begin{equation}
                        \label{eqGalerUpdateRuleNewtonMethod}
                        \mBF u_{\delta,n+1}^\boldmu
                        =\mBF u_{\delta,n}^\boldmu
                        +\delta\mBF u_n^\boldmu
                    \end{equation}
                    %
                    where now every term is are vectors carrying the weights associated to the basis functions;

                \item $\mBF A$ is the stiffness matrix ($A_{ij}\equiv a(\varphi_j,\varphi_i)$), $\mBF b$ is $\mBF b_j\equiv b(1,\varphi_j)$ while $\mBF b^{\mu_1}_n$ is $(\mBF b^{\mu_1}_n)_j\equiv b(e^{\mu_1u^\boldmu_{\delta,n}},\varphi_j)$;
                
                \item $\mBF A$, $\mBF b$ and $\mBF g$ are always affine w.r.t. $\boldmu$;
                
                \item $\mBF B_n^{\mu_1}$ and $\mBF b_n^{\mu_1}$ are always \textit{not} affine w.r.t. $\mu_1$;
                
                \item $\mBF g^\boldmu_n$ is affine w.r.t. $\mu_0$ if $g\equiv g_1$ but \textit{not} if $g\equiv g_2$.
            \end{itemize}

        \subsection{Proper Orthogonal Decomposition (POD)}\label{secPOD}

            POD is strategy that aims to reduce computational costs by solving the system \eqref{eqGalerIterLS} in a reduced space $V_\rho\subseteq V_\delta$ of dimension $\Nrho\ll\Ndelta$: assume to know the basis function $\{\xi_k\}_{k=1}^\Nrho\in V_\delta$ of such reduced space (i.e. $V_\rho=\spanOp\{\xi_k\}_{k=1}^\Nrho$); then it's possible to express each function w.r.t. $\{\varphi_k\}_{k=1}^\Ndelta$, obtaining the basis matrix 
            $$
                \ReducedBase=[\boldxi^1\cdots\boldxi^\Nrho]\in\mBB R^{\Ndelta\times\Nrho};
            $$
            given a vector $v\in\mBB R^\Nrho$, one has $\ReducedBase v\in\mBB R^\Ndelta$, that is $\ReducedBase$ can be seen as a projection from $V_\rho$ to $V_\delta$ and the same can be said for the inverse projection with $\ReducedBase^\top$; thus, defining $\delta\hat{\mBF u}^\boldmu_n$ such that $\delta\mBF u^\boldmu_n=\ReducedBase\delta\hat{\mBF u}^\boldmu_n$, \eqref{eqGalerIterLS} becomes:
            \begin{equation}
                \label{eqReducedGalerIterLS}
                \hat{\mBF S}^\boldmu_n\delta\hat{\mBF u}^\boldmu_n
                =\hat{\mBF y}^\boldmu_n,
            \end{equation}
            %
            in which $\hat{\mBF S}^\boldmu_n\equiv\ReducedBase^\top\mBF S^\boldmu_n\ReducedBase$ and $\hat{\mBF y}^\boldmu_n\equiv\ReducedBase^\top\mBF y^\boldmu_n$. The strategy can be summarized using the following commutative diagram: 
            $$
                \begin{tikzcd}[
                    row sep=1cm,
                    column sep=2cm,
                ]
                    V_\delta
                    \arrow[r,"\mBF S^\boldmu_n\delta\mBF u^\boldmu_n=\mBF y^\boldmu_n"] 
                    \arrow[d,"\ReducedBase^\top"']&
                    V_\delta\arrow[d,"\ReducedBase",leftarrow]\\
                    V_\rho
                    \arrow[r,"\hat{\mBF S}^\boldmu_n\delta\hat{\mBF u}^\boldmu_n=\hat{\mBF y}^\boldmu_n"]&
                    V_\rho
                \end{tikzcd},
            $$
            where one hopes the reduced path is less computationally demanding than the full one. Of course, as mentioned before, both the iterative and non-affine nature of \eqref{eqGalerIterLS} hinders this chance, though only numerical studies in \S{} \ref{secNumericalStudies} can confirm if that's truly the case.

            To find $\ReducedBase$ consider $M\ll\Ndelta$, but large enough so that $M\gg\Nrho$, FOM solutions $\mBF u^n_\delta\equiv\mBF u^{\boldmu_n}_\delta$, $1\le n\le M$, of \refVarGalerNLEllPDE{} with uniformly varying parameters $\boldmu$, called snapshots, and define
            $$
                \mBB U\equiv
                \begin{bmatrix}
                    \mBF u^1_\delta&
                    \cdots&
                    \mBF u^M_\delta
                \end{bmatrix}^\top
                \in\mBB R^{M\times\Ndelta}
            $$
            as the matrix whose rows are the snapshots.
            
            Much of the information inside $\mBB U$ is redundant, so one then needs to compress it via the correlation matrix
            %
            \begin{equation}
                \label{eqCovarianceMatrixDefinition}
                \mBF C\equiv\mBB U\mBF A\mBB U^\top\in\mBB R^{M\times M};
            \end{equation}
            %
            indeed, by solving the eigenvalue problem
            %
            \begin{equation}
                \label{eqCovarianceMatrixEigenProblem}
                \mBF C\boldomega^i=\lambda_i\boldomega^i,\quad1\le i\le M,
            \end{equation}
            %
            it holds
            %
            \begin{equation}
                \label{eqErrorCovarianceEigenVectors}
                \text{err}^2_{M,\rho}
                =\sum_{i=1}^M
                \Vert u_\delta^i-P_\rho(u_\delta^i)\Vert_V^2=
                \sum_{i=\Nrho+1}^M\lambda_i,
            \end{equation}
            %
            where $P_\rho:V_\delta\to V_\rho$ is the projection operator to the reduced space. In essence \eqref{eqErrorCovarianceEigenVectors} says that the most important eigenvectors, \ie{} those that retain all significant variance, are those with the largest eigenvalues; and since $\mBF C$ is real, symmetric and positive semi-definite, we can order the all-positive eigenvalues as $$\lambda_1\ge\dots\ge\lambda_M\ge0$$ and take only the first $\Nrho$ so that $\text{err}^2_{M,\rho}\approx0$.

            However, the eigenvectors $\{\boldomega^k\}_{k=1}^\Nrho$ are still in $\mBB R^M$ while $\{\boldxi^k\}_{k=1}^\Nrho$ lie in $\mBB R^\Ndelta$; to transform them accordingly it is sufficient to notice how \eqref{eqCovarianceMatrixEigenProblem} changes with \eqref{eqCovarianceMatrixDefinition}:
            %
            \begin{equation}
                \label{eqTransformedCovarianceMatrixEigenProblem}
                \newcommand\implica{\;\;\implies\;\;}
                \begin{aligned}
                    \eqref{eqCovarianceMatrixEigenProblem}&\implica
                    \mBB U\mBF A\mBB U^\top\boldomega^i=\lambda_i\boldomega^i,
                    \\&\implica
                    \mBF A\mBB U^\top\boldomega^i=\lambda_i\mBB U^\top\boldomega^i,
                    \\&\implica
                    \mBF A\boldchi^i=\lambda_i\boldchi^i,
                    \quad1\le i\le M,
                \end{aligned}
            \end{equation}
            %
            where $\boldchi^i=\mBB U^\top\boldomega^i$, $1\le i\le N_\rho$, is the transformation sought which also proves $\boldchi^i$ are the eigenvectors of the stiffness matrix
            \footnote{The stiffness matrix shares many properties of $\mBF C$: by definition ($A_{ij}\equiv a(\varphi_j,\varphi_i)$) it's real, symmetric and due to the coercivity of $a$ ($a(v,v)\ge\alpha\Vert v\Vert^2_V\;\;\;\forall v\in V$) it's also definite positive.}.
            
            The last step comes from linear algebra: in fact, we know thnaks to the spectral theorem that both ${\boldchi^k}_{k=1}^\Nrho$ and $\{\boldomega^k\}_{k=1}^\Nrho$ are only orthogonal to one another\cite[\S 6.3 p. 286]{giordano2017geometria}; hence the reduced basis vectors are obtained by simply normalizing them:
            $$
                \boldxi^i\equiv\frac{\boldchi^i}{\Vert\boldchi^i\Vert_V}.
            $$
            
        \subsection{Physics-Informed Neural Networks (PINN)}\label{secPINN}
    
            In non-linear frameworks, to overcome some POD method's limitations and to exploit the advantages of the division in online and offline phases, the use of the neural networks has opened new opportunities for solving PDEs. In particular, PINN is a flexible approach that directly uses the equations into the training process of a neural network, it incorporates physical laws given by the differential equations into the loss functions, by considering the residual of the PDE, to take advantages in the learning process. This allows PINNs to deal with complex and parametric problems in a data efficient and mesh free way.

            The main idea is to incorporate into the loss function both the information obtained from the data and the physical model. The physics-informed term is like a prior knowledge that provides more information to the PINN without the need for more measurements, but only by evaluating the residual of the PDE at additional points in the domain. In our case, we want to approximate the solution of the PDE by considering two contributions to the loss function: one related to boundary information and the other to residual. 
            
            The general structure of the parametric PINN method can be summarized as follows:
            %
            \begin{equation}
                \label{eq:NN}
                \begin{aligned}
                    \mathcal{F}_{NN} \colon \quad & \Omega \times \mathcal{P} \to \mathbb{R}^{N_\delta}, \\
                    & (\mathbf x,\boldmu) \mapsto \hat{\mathbf u}^\boldmu(\mathbf x) \approx \mathbf u^\boldmu(\mathbf x).
                \end{aligned}
            \end{equation}
            %
            where $ \mathbf u^\boldmu(\mathbf x)$ is the vector of coefficients associated with the Galerkin framework, whose entries correspond to the weights of the basis functions. 
            The neural network is a fully-connected feed-forward neural network with $n_{\mathrm{HL}}$ hidden layers, each of width $n_{\mathrm H}$.

            The training process is given by the minimization of the loss function define as:
            %
            \begin{equation}
                \MSE=\MSE^\boldmu_b+\lambda \MSE^\boldmu_p
            \end{equation}
            %
            where $\MSE$ is the Mean Squared Error and we compute it for the boundary points using $\MSE^\boldmu_b$, this enforces the condition $u^\boldmu=0$, while for the inner points the residual loss $\MSE^\boldmu_p$. Finally, $\lambda$ is the hyperparameter that balances the trade-off between the two contributions, to be set before the network training.

            
            In particular, by considering our problem, the errors are given by:
            %
            \begin{equation*}
                \MSE^\boldmu_b=\frac{1}{N_b}\sum_{k=1}^{N_b}\left|\hat u^\boldmu(\mathbf x^b_k;\theta)-u^\boldmu(\mathbf x^b_k)\right|^2,
            \end{equation*} 
            %       
            \begin{equation*}
               \MSE^{\boldmu}_{p}=\frac{1}{N_p}\sum_{k=1}^{N_p} \left| -\Delta \hat u^{\boldmu}(\mathbf x^p_k)+h^\boldmu(\hat u^{\boldmu}(\mathbf x^p_k)) -g^{\boldmu}(\mathbf x^p_k)-\frac{\mu_0}{\mu_1}\right|^2
            \end{equation*}
            %
            for $x^{b}_k \in \partial \Omega$ and $x^p_k \in \Omega$ are the selected points in which we compute the loss and $N_b$, $N_p$ denote respectively the number of boundary and interior points. This is a mesh-free approach, so we need to randomly select the boundary points and interior points by fixing the number of them, which is needed to build the network.  
            The hyperparameter $\lambda$ balances the relative importance of the two contributions.  
            Rather than fixing $\lambda$ a priori, we dynamically adapt its value during the training. It is updated at each iteration depending on the relative magnitude of $\MSE^\boldmu_b$ and $\MSE^\boldmu_p$. In practice, if one loss term dominates the other, $\lambda$ is modified so as to increase the weight of the weaker contribution, in this way we can balance in a better way the losses and we can be sure that the network is not optimized by relying too much neither over the boundary part nor over the inner part. In this way, the overall stability of the training process is improved.
    
        \subsection{POD-NN hybrid approach}\label{secPODNN}
    
            By combining the previous two methods, we obtain a promising paradigm: the POD-NN hybrid approach. This leverages the dimensionality reduction given by POD to obtain a compact representation of the solution space while using a neural network to learn the parametric dependence of the reduced coefficients. This strategy maintains the accuracy of POD while exploiting the benefits of neural networks, providing a computationally efficient alternative to both the MOR and PINN approaches.

            POD-NN is a strategy that provides a fast tool to deal with non affine structures. It avoids the online projection by exploiting the direct projection of snapshots onto the reduced space built with the POD method, and then using them to create and train a neural network. 
            
            The general structure of the network can be summarized as follows:
            %
            \begin{equation}
                \label{eq:podNN}
                \begin{aligned}
                    \pi_{NN} \colon \quad & \mathcal{P} \to \mBB{R}^{N}, \\
                    &\boldmu^* \mapsto \boldsymbol u^{NN}_{N}(\boldmu^*) 
                \end{aligned}
            \end{equation}
            %
            where $N$ is the dimension of the reduced space obtained by applying the POD. The inputs are the parameters of the training set, while the output is the Galerkin projection of the snapshots. 
            For the output, which has the form $\boldsymbol u^{NN}_N$, the following relation holds:
            %
            \begin{equation*}
                \mBB B \boldsymbol u^{NN}_{N} (\boldsymbol \mu) \approx \mBB P^{\boldsymbol \mu} \boldsymbol u_{{\delta}}(\boldsymbol \mu)
            \end{equation*}
            %
            where $\mBB B \in \mBB{R}^{N_\delta \times N}$ is the basis functions matrix of the reduced space, and $\mBB P^{\boldsymbol \mu}=\mBB B \mBB P_{\delta \to N} \in \mBB{R}^{N_\delta \times N_\delta}$ is composed by the projection operator, which leads to the best approximation of $u^\boldmu_\delta$ in the reduced space, pre-multiplied by the basis function matrix. This allows us to compute $\boldsymbol u^{NN}_{N} (\boldsymbol \mu)$ for each snapshot without solving the reduced system, obtaining the best approximation to $\boldsymbol u_{\delta}$ because we have a direct connection between $ \boldsymbol u^\boldmu_\delta$ and $\boldsymbol u^{NN}_{N} (\boldsymbol \mu)$.
            To train the network, we minimize the loss given by the distance between the exact $u_N$ solution and the network approximation:
            %
            \begin{equation}
                \label{eq:PODNNLoss}
                \MSE=\frac{1}{N_{snap}} \sum_{\boldsymbol \mu \in \mathcal{P}_{train}} 
                \lvert \lvert \boldsymbol u^{NN}_N(\boldsymbol \mu) - \boldsymbol u_N(\boldsymbol \mu) \rvert \rvert^2 ,
            \end{equation}
            %
            where $N_{snap}$ is the number of data in the training set. 
            
            To sum up, we can compute offline the snapshots collection, apply the POD strategy and train the network. Then, we evaluate the network during the online phase. In this way, we exploit the POD properties and the online phase becomes very fast. In contrast, the offline phase is longer and the accuracy is more difficult to control.  

    %
    
    \section{Numerical studies}\label{secNumericalStudies}
    
        In this section the reduced solution seen in § \ref{secNumericalStudies} are compared in terms of computational costs and accuracy with respect to the FOM one.

        \textbf{
            Introduzione con 
            \begin{itemize}
                \item 'set up' di simulazione
                \item metriche usate per valutare risultati
                \item confronti tra gli esperimenti 
            \end{itemize}
            Esperimenti nello specifico: confronti risultati con grafici
        }
        
        \subsection{Source term independent of parameters}\label{secIndependentCase}
        
        \subsection{Source term dependent on parameters}\label{secDependentCase}

    %
    
    \section{Conclusions}\label{seConclusion}

    %

    \bibliographystyle{IEEEtran}
    \bibliography{../References/Bibliography}
    \nocite{*}

\end{document}